% ============================================================
%  HCMUT-DEE Thesis Kit v2.0.4 — Tóm tắt tiếng Việt
%  Copyright (c) 2026 Nguyễn Trọng Thắng & TS. Nguyễn Phúc Khải
%  License: MIT
% ============================================================


% ═══════════════════════════════════════════════════════════════
%  HUONG DAN: Viet tom tat do an cua ban o day.
%  Noi dung tom tat thuong bao gom:
%    - Van de nghien cuu (1-2 cau)
%    - Phuong phap thuc hien (2-3 cau)
%    - Ket qua chinh (2-3 cau)
%    - Y nghia va dong gop (1-2 cau)
%
%  Do dai khuyen nghi: 1 trang (khoang 200-300 tu).
%  Khong dung viet tat (tru khi giai thich lan dau).
%  Khong dung trich dan tai lieu tham khao.
% ═══════════════════════════════════════════════════════════════

\chapter*{TÓM TẮT DỰ ÁN}
\addcontentsline{toc}{chapter}{Tóm tắt dự án}

Việc soạn thảo báo cáo kỹ thuật và đồ án tốt nghiệp luôn là một thách thức lớn đối với sinh viên ngành kỹ thuật do các yêu cầu khắt khe về định dạng văn bản học thuật. Các phương pháp soạn thảo truyền thống như Microsoft Word thường bộc lộ nhiều hạn chế về tính đồng bộ, lỗi định dạng khi thay đổi môi trường làm việc, và khó khăn trong việc quản lý mã nguồn, hình ảnh hay tài liệu tham khảo. Để giải quyết triệt để rào cản này, dự án đã xây dựng và chuẩn hóa bộ công cụ **HCMUT-DEE Thesis Kit** – giải pháp template LaTeX tối ưu được thiết kế chuyên biệt cho sinh viên Khoa Điện - Điện tử, Trường Đại học Bách Khoa - ĐHQG-HCM.

Thuyết minh này đóng vai trò là tài liệu hướng dẫn sử dụng và là mẫu cấu trúc chính thức của bộ template. Bộ công cụ được phát triển dựa trên cấu trúc lớp tài liệu (`hcmut-dee.cls`) kế thừa từ nền tảng `report` tiêu chuẩn, tích hợp sẵn các gói lệnh thiết yếu để xử lý toán tử phức tạp, tự động tạo trang bìa (\texttt{\textbackslash makecover}), trang nhiệm vụ (\texttt{\textbackslash makeassignment}) và tờ chấm (\texttt{\textbackslash makeevaluation}). Đặc biệt, template cung cấp các cấu hình định dạng mã nguồn nâng cao dành cho Python, MATLAB, C/C++ và ngôn ngữ PLC (Structured Text), đáp ứng trọn vẹn nhu cầu đặc thù của các phân ngành Hệ thống điện, Điện tử, Viễn thông và Tự động hóa.

Kết quả triển khai thực tế cho thấy bộ template giúp sinh viên giảm thiểu hơn 80\% thời gian dành cho việc định dạng văn bản, loại bỏ hoàn toàn các lỗi tràn lề bảng biểu hay trích dẫn sai quy chuẩn, đồng thời tăng tính chuyên nghiệp và thẩm mỹ của tài liệu. Dự án hướng tới tinh thần chia sẻ và truyền thừa tri thức qua các thế hệ sinh viên, đồng thời thúc đẩy việc ứng dụng LaTeX trong nghiên cứu khoa học kỹ thuật tại nhà trường.

\vspace{0.5cm}
\noindent \textbf{Từ khóa:} LaTeX Template, Báo cáo kỹ thuật, Đồ án tốt nghiệp, Khoa Điện - Điện tử, ĐH Bách Khoa, Kế thừa tri thức.
